% AI-produced by GPT-6 Astra (OpenAI, ChatGPT Pro mode). No priority claim.
\subsection{The exact component-injectivity threshold}

The obstruction left by
Corollary~\ref{C:equalizer-component-reduction} can be stated
entirely in terms of the old component maps. In particular one
must control long component paths as well as component cycles.

\begin{proposition}[Proved: injectivity in a set necklace]
\label{C:set-necklace-injectivity}
Let $f,g:S\rightrightarrows T$ be injections of finite sets, and
let $s=4k+1$, $k\geq1$. Form the set coequalizer of the cyclic
diagram with $s$ copies of $T$ and neighboring legs
\[
 (\lambda_i,\mu_i)=
 \begin{cases}
 (g,f),&0\leq i<2k,\\
 (g,g),&i=2k,\\
 (f,g),&2k<i<4k,\\
 (f,f),&i=4k.
 \end{cases}
\]
Write $j_i:T\to Z$ for its copy maps. Then all the maps $j_i$
are injective if and only if the partial relation
$f(x)\mapsto g(x)$ has no nontrivial directed cycle and each
nonconstant directed path has at most $2k+1$ vertices.
\end{proposition}

\begin{proof}
Because $f,g$ are injections, their partial transition relation
decomposes into fixed points, directed paths, and nontrivial
directed cycles. Points outside both images are private. Every
neighbor identification stays within one of these pieces, so
different pieces can be treated separately. A fixed point is
shared by all copies, while private points remain separate.

On a path of $l$ vertices, number them $0,\ldots,l-1$ so that
the source transitions are $t\mapsto t+1$ for
$0\leq t\leq l-2$. Suppose $l\leq2k+1$. The same four-branch
formula as in the repaired construction defines a cocone into
$\{0,\ldots,4k\}$:
\[
 b_i(t)=
 \begin{cases}
 t+i,&i\leq2k,\ t+i\leq2k,\\
 t+i+2k-l+1,&i\leq2k,\ t+i>2k,\\
 t-i+6k-l+2,&i>2k,\ t-i+6k-l+2>2k,\\
 t-i+4k,&i>2k,\ t-i+6k-l+2\leq2k.
 \end{cases}
\]
The first-arm maps are increasing: the jump between their two
branches is nonnegative in addition to the unit increase,
because $2k-l+1\geq0$. The second-arm maps are also increasing
in $t$: the change in their branch offsets is
$2k-l+2\geq1$. The displayed values lie between $0$ and $4k$.
Thus each $b_i$ is injective.

For $i<2k$, the equality $b_i(t+1)=b_{i+1}(t)$ follows by
substitution. At the upper junction,
\[
       b_{2k}(t+1)=t+4k-l+2=b_{2k+1}(t+1);
\]
the last value uses its upper branch since
$4k-l+2\geq2k+1$. On the second arm,
$b_i(t)=b_{i+1}(t+1)$. Finally
$b_{4k}(t)=t=b_0(t)$ for $0\leq t\leq l-2$.
These are all the required neighbor identities. Since the
coequalizer has a cocone whose restriction to each copy is
injective, its own copy maps cannot identify distinct points
within a copy. Separate tags for the different path pieces,
and the fixed and private pieces, prove sufficiency for the
whole diagram. This argument only needs $l\leq2k+1$, which is
weaker than the usual sufficient hypothesis $l\leq k$.

For necessity, first suppose a path has at least $2k+2$ vertices.
Start at position $2k+1$ in copy zero. Moving forward across the
first $2k$ neighboring identifications decreases the position
by one each time, reaching position $1$ in copy $2k$.
The $g,g$ junction keeps position $1$. The following $2k-1$
identifications increase the position by one each time, reaching
position $2k$ in copy $4k$. The $f,f$ junction returns this to
position $2k$ in copy zero. All used positions are permitted by
the assumed path size. Hence the coequalizer identifies the two
distinct points $2k+1$ and $2k$ of copy zero.

On a directed cycle, the same circuit of identifications is
always available, with positions interpreted modulo its length.
The net change around the index necklace is minus one: the first
arm contributes $-2k$, the second contributes $2k-1$, and the two
junctions contribute zero. Every point of that cycle is thus
identified with its predecessor in the same copy. A nontrivial
cycle makes the copy map noninjective. This proves necessity.
\end{proof}

\begin{corollary}[Proved: the first possible component failure]
\label{C:component-necklace-obligation}
Assume every arrow of an old repaired C.1 stage is injective on
$H(X)=\pi_0G_X$. For an attachment to $f,g:S\rightrightarrows T$
with $s=4k+1$, every new generating map is injective on $H$ if
and only if the partial relation
\[
                           H(f)(c)\mapsto H(g)(c)
\]
has no nontrivial directed cycle and every path has at most
$2k+1$ vertices. Under these conditions every new arrow is
injective on $H$.

Consequently the first failure of component injectivity in a
C.1 construction, if there is one, must occur at an attachment
whose old component relation has a nontrivial cycle or a path
with at least
\[
                        2k+2=\frac{s+3}{2}
                                   \geq\frac{15n+3}{2}
\]
vertices. Theorem~\ref{C:universal-components} excludes both
possibilities for the actual generated arrows: their component
relations are acyclic and their path sizes are at most
$L_U(f,g)+1\leq n+1\leq k$.
\end{corollary}

\begin{proof}
The repaired graph attachment is the strict graph colimit of its
necklace diagram: vertices have the set-coequalizer description,
and $(c_m)$ routes the edge fibers through the same neighboring
overlaps, as in Proposition~\ref{C:fixed-path-kernel}.
Passing to graph components preserves this colimit. Thus
$H(Z)$ is precisely the set necklace for $H(f),H(g)$, with the
same $s$ and the same four kinds of neighboring legs. The old
component maps are injections by hypothesis, so
Proposition~\ref{C:set-necklace-injectivity} applies. Every
new arrow factors through a generating copy, proving the
assertion about all new arrows. Sum targets preserve component
injectivity directly. Finally C.1 gives $s\geq15n$, as in
Lemma~\ref{C:birth-cover-helly}.
\end{proof}

\paragraph{Verification and scope.}
The retained program
\texttt{akr\_component\_necklace\_threshold.py} constructs the
set coequalizer directly by union--find for $k=1,2,3,5,11$ and
path sizes from two through $4k+4$. In each case the first
failure is exactly $2k+2$. This is an independent finite check
of the threshold; the proof above is uniform in $k$ and does
not rely on those tests. The proposition and its corollary
isolate the exact set-level threshold. The separate simultaneous
induction of Theorem~\ref{C:universal-components} supplies the
all-stage bound on the actual component-transition relations.
